\chapter*{Presentacion}
\addcontentsline{toc}{section}{Presentacion}

Este cuaderno nace del corpus real disponible en \texttt{sources/}. A partir de esas relaciones
de ejercicios se ha construido una taxonomia reproducible, una matriz de cobertura y una
redaccion didactica que convierte el material disperso en un libro de repaso con progresion.

La obra queda organizada en tres capas complementarias:
\begin{itemize}
  \item \textbf{Capitulos nucleares \texttt{C01-C09}.} Recogen el temario matematico respaldado
  por fuentes en la carpeta actual.
  \item \textbf{Capitulo original de estadistica y probabilidad.} Completa el sentido
  estocastico con datos bidimensionales, regresion, muestreo, recuento y Bayes. Al no existir un
  PDF fuente de este bloque, se identifica expresamente como contenido editorial original.
  \item \textbf{Repaso acumulativo.} Recombina tecnicas de varios bloques para entrenar eleccion
  de metodo, deteccion de errores y simulacros con tiempo orientativo.
  \item \textbf{Tres salidas coordinadas.} La edicion del alumnado conserva teoria, guiados y
  respuestas breves; la del profesorado anade las soluciones completas; y el documento de
  respuestas sirve como consulta rapida.
\end{itemize}

La trazabilidad documental de \texttt{C10} no puede vincularse a una relacion fuente porque no
existe ese PDF en la carpeta de trabajo. Si se incorpora en el futuro, habra que rehacer la
auditoria y la matriz de cobertura para contrastar el capitulo original con el nuevo corpus.
